% !PRINTED 457
\weberpage{457}
The straight line $x_1 = 0$ cuts the curve in three coincident
points at $x_3 = 0$ and in a fourth point $x_2 = 0$ separated from
them. The vertices of the coordinate triangle are therefore points of
inflection of the curve, and the sides are the inflectional tangents.
If we denote the sides of the coordinate triangle by 1, 2, 3 and the
opposite vertices by the same figures, then side 1 is the inflectional
tangent at point 2, side 2 the inflectional tangent at point 3, and
side 3 the inflectional tangent at point 1.

We now investigate the position of the poles and axes with respect to
the fundamental curve. Since the vertices of the coordinate triangle
belong to the sevenfold poles, we conclude first of all
\gesperrt{that all sevenfold poles $P_7$ lie on the fundamental
curve.}

They are, like the vertices themselves, all \gesperrt{points of inflection} of
the curve, and these arrange themselves into eight inflection-point
triangles.

\gesperrt{The threefold poles} $P_3$ \gesperrt{likewise lie on the fundamental curve;
for if one sets}
\[
x_2 = \varrho\,x_1, \quad x_3 = \varrho^2 x_1
\]
or
\[
x_2 = \varrho^2 x_1, \quad x_3 = \varrho\,x_1,
\]
in which $\varrho$ is a cubic root of unity, then
$f(x_1, x_2, x_3)$ reduces to
\[
1 + \varrho + \varrho^2 = 0.
\]
To find the points of intersection of their connecting line
\[
T = x_1 + x_2 + x_3 = 0
\]
with the fundamental curve, one sets $x_3 = -x_1 - x_2$, whereby
$f(x_1, x_2, x_3) = 0$ passes into
\[
x_1^3(x_1+x_2) - x_2^3 x_1 + (x_1+x_2)^3 x_2 = (x_1^2+x_1x_2+x_2^2)^2 = 0
\]
Since the left side is a square, the four points of intersection
coincide twice in pairs, and it follows that the line $T$ is a
double tangent of the fundamental curve.

\gesperrt{The 28 lines $T$ are double tangents of the fundamental
curve; their points of contact are the 56 poles $P_3$.}

From this it follows incidentally that the points $P_7$ do not lie
on the lines $T$, since a double tangent can have no further point
of intersection with the curve besides the points of contact.

\gesperrt{The six- and eightfold poles do not lie on the fundamental
curve.}
